\subsection{Modelo de detecção}

A frente de aprendizado de máquina seguiu um caminho deliberadamente simples,
um passo de cada vez. Primeiro mostrou que o RSSI legítimo pode ser previsto
melhor que pela regra mais simples possível, e só depois mediu a detecção. As
atividades foram as seguintes, nesta ordem.

\begin{itemize}
    \item \textbf{Preparação da série.} limpeza do \textit{dataset} Tour Perret
    com nove regras contadas e recorte de uma série de RSSI de um nó sensor num
    \textit{gateway}, descritos na Seção~\ref{subsec:ml-limpeza}.

    \item \textbf{Análise da série.} ACF, PACF e teste KPSS no treino, para
    decidir a diferenciação.

    \item \textbf{Previsão.} ARIMA na série inteira, comparado ao
    \textit{random walk}, e depois um ARIMA por canal de frequência, com a
    ordem variada um parâmetro de cada vez na validação, avaliados pelo RMSE,
    pelo MAE e pela mediana do erro absoluto.

    \item \textbf{Quantidade de dados.} treinos de 6~h a 142 dias, para medir
    quanto um nó na borda precisa coletar antes de prever.

    \item \textbf{Detecção.} limiar por canal calibrado na validação, desvios
    sintéticos injetados em 300 posições do teste e comparação com um detector
    mais simples, baseado no \textit{random walk} por canal.

    \item \textbf{Reprodutibilidade.} cada número gravado por \textit{script} e
    ligado ao arquivo de origem por resumo SHA-256, com reexecução completa numa
    cópia isolada do projeto.
\end{itemize}

Os resultados estão na Seção~\ref{sec:resultados_ml}.
